% ==========================================================================
%  Chapter 84 -- Execution record for multiscale stabilization v1
% ==========================================================================

\chapter{Execution Record for Multiscale Stabilization v1}
\label{ch:multiscale-stabilization-execution}

\section{Intent and Verification Protocol}
\label{sec:ms-stab-protocol}

This chapter records the actual execution of the multiscale stabilization plan
after the publication-oriented redesign documented in
Chapter~\ref{ch:multiscale-publication-audit}.  The guiding rule for this
stabilization cycle was strict: neither the previous code, nor the existing
tests, nor the earlier documentation were accepted as ground truth without a
fresh audit against the implemented path.

The execution protocol used in this cycle was:
\begin{enumerate}
  \item audit the current normative path directly in code;
  \item reproduce uncertainty or weakness with a dedicated test or diagnostic;
  \item implement the correction or hardening;
  \item rerun a focused regression gate;
  \item document both the improvement and the residual debt in the same
    \LaTeX{} document.
\end{enumerate}

The main files audited and modified during this cycle were:
\begin{itemize}
  \item \texttt{src/analysis/MultiscaleTypes.hh};
  \item \texttt{src/analysis/MultiscaleAnalysis.hh};
  \item \texttt{src/reconstruction/NonlinearSubModelEvolver.hh};
  \item \texttt{src/reconstruction/LocalModelAdapter.hh};
  \item \texttt{src/reconstruction/HomogenisationStrategy.hh};
  \item \texttt{tests/test\_multiscale\_analysis\_api.cpp};
  \item \texttt{tests/test\_evolver\_advanced.cpp};
  \item \texttt{scripts/ci\_multiscale\_stabilization.ps1};
  \item \texttt{.github/workflows/multiscale-stability.yml}.
\end{itemize}

\section{Phase 0: Truth Baseline}
\label{sec:ms-stab-phase0}

Table~\ref{tab:ms-stab-truth-baseline} summarises the baseline classification
used before any code change was accepted.

\begin{table}[htb]
\centering
\small
\begin{tabular}{p{0.27\textwidth}p{0.13\textwidth}p{0.43\textwidth}}
\toprule
Audited item & Status & Decision in this stabilization cycle \\
\midrule
Typed multiscale orchestration path &
verified &
Kept as the normative public path. \\
Type-erased \texttt{LocalModelHandle} path &
experimental &
Left in code, but explicitly removed from the normative public narrative until
exact erased checkpoint/restore semantics are implemented. \\
\texttt{HomogenisationStrategy} route &
legacy/internal &
Retained only as a comparison utility; no longer described as the production
FE\textsuperscript{2} extension mechanism. \\
Local tangent extraction in
\texttt{NonlinearSubModelEvolver} &
fragile &
Replaced by adaptive finite differences with central preference, one-sided
fallback, and explicit diagnostics. \\
Macro rollback on trial-step failure inside FE\textsuperscript{2} &
incorrect &
Explicit checkpoint restoration added to the failure path. \\
Physical regression matrix for pure load modes &
missing &
Added as dedicated solver-level regression tests. \\
Visible CI gate for the multiscale release surface &
missing &
Added as a documented local script and a repository workflow. \\
\bottomrule
\end{tabular}
\caption{Truth baseline for the stabilization cycle.}
\label{tab:ms-stab-truth-baseline}
\end{table}

\section{Phase 1: Numerical Hardening of the Local Response}
\label{sec:ms-stab-phase1}

\subsection{Initial hypothesis}

The main numerical fragility of the public multiscale path was the local
tangent extraction: a one-sided perturbation with a fixed user-provided
\(h_{\mathrm{pert}}\), silent zero columns on failed perturbations, and hidden
post-processing by symmetry projection and diagonal clamping.

\subsection{Observed evidence before the change}

The audit of \texttt{NonlinearSubModelEvolver.hh} showed that the previous
implementation:
\begin{itemize}
  \item used only the \(+\) perturbation;
  \item used the same perturbation size for every generalized component;
  \item left the tangent column at zero if the perturbation solve diverged;
  \item restored displacement vectors but did not explicitly resynchronise the
    model state vector in the tangent loop;
  \item reported only the final tangent, thereby hiding which columns were
    valid and whether regularisation had actually modified the matrix.
\end{itemize}

\subsection{Implemented change}

The production operator remains the boundary-reaction operator, but the
stabilization no longer treats finite differences as the primary tangent path.
The current normative implementation first assembles the full local tangent
operator at the converged micro state, including the embedded-rebar penalty
terms, partitions it into free and constrained blocks, and condenses it
against the boundary kinematics.  With constrained increments
\(\delta \mathbf{u}_c = \mathbf{T}\,\delta \mathbf{q}\), the returned section
tangent is
\begin{equation}
  \mathbf{D}_{\mathrm{hom}}
  =
  \mathbf{T}^{\mathsf T}
  \mathbf{K}_{\mathrm{eff}}
  \mathbf{T},
  \qquad
  \mathbf{K}_{\mathrm{eff}}
  =
  \mathbf{K}_{cc}
  -
  \mathbf{K}_{cf}\mathbf{K}_{ff}^{-1}\mathbf{K}_{fc}.
\end{equation}
This is the first version of the multiscale operator in the repository that is
explicitly built from the linearised micro equilibrium equations rather than
from perturb-and-observe solves.

Adaptive finite differences are retained only as a documented fallback.  When
the condensed operator cannot be formed---for example because a mixed element
does not expose explicit local tangent data or because the free block is
numerically singular---the implementation falls back to the adaptive scheme.
For each generalized component \(e_j\), the fallback uses
\begin{equation}
  h_j
  =
  \frac{\max\!\left(h_{\mathrm{pert}} L,\;
  10^{-3}\max(|q_j|, h_{\mathrm{pert}}L)\right)}{L},
\end{equation}
where \(L\) is the macro beam characteristic length and \(q_j\) is the
corresponding relative end degree of freedom.  The preferred tangent estimate
is
\begin{equation}
  \mathbf{D}_{\mathrm{hom}}^{(:,j)}
  \approx
  \frac{\mathbf{s}(\mathbf{e}+h_j \mathbf{e}_j)
       -\mathbf{s}(\mathbf{e}-h_j \mathbf{e}_j)}{2h_j},
\end{equation}
and the implementation falls back to one-sided differentiation only when one of
the two perturbation solves fails.

The public response object now records:
\begin{itemize}
  \item the tangent scheme used;
  \item the perturbation size used for each column;
  \item whether each column is valid;
  \item whether each column used the central stencil;
  \item the number of failed perturbation solves;
  \item whether regularisation actually modified the returned tangent.
\end{itemize}

\subsection{Resulting semantics}

The stabilization deliberately stopped calling every non-perfect response
``successful''.  The public status contract is now:
\begin{description}
  \item[\texttt{Ok}] the condensed operator was assembled successfully, all
    tangent columns are valid, and no actual regularisation was applied;
  \item[\texttt{Degraded}] a usable response exists, but at least one
    fallback to finite differences or actual tangent regularisation occurred;
  \item[\texttt{SolveFailed}] the response is not usable as a section operator;
  \item[\texttt{NotReady}/\texttt{InvalidOperator}] explicit non-solution
    states.
\end{description}

\subsection{Internal architectural hardening}

The numerical hardening was paired with an internal architectural split of the
local solver path.  Instead of leaving boundary conditions, persistent state
management, and boundary-reaction homogenisation interleaved inside
\texttt{NonlinearSubModelEvolver}, the cycle extracted three dedicated
components:
\begin{itemize}
  \item \texttt{PersistentLocalStateOps}: synchronisation, commit, and revert
    of the persistent continuum state;
  \item \texttt{LocalBoundaryConditionApplicator}: update of section-end
    kinematics and writing of the imposed boundary vector;
  \item \texttt{BoundaryReactionHomogenizer}: section forces from reactions,
    condensed tangent extraction, finite-difference fallback, diagnostics, and
    response classification.
\end{itemize}

This split was intentionally implemented as lightweight header-only helpers
composed through pointers and templates, so the hot path still avoids virtual
dispatch and heap churn.  Architecturally, this matters because it creates a
single explicit seam where the current dense condensation / finite-difference
fallback pair can later be upgraded again without changing the public
\texttt{SectionHomogenizedResponse} contract or the typed local model API.

The same cycle also extended \texttt{FEM\_Element} with the explicit local
linearisation API already present in \texttt{StructuralElement}.  This closes
an architectural gap in the mixed-element path: the condensed local operator
now sees the element DOF map and the element tangent even when the micro model
is stored behind \texttt{MultiElementPolicy} type erasure.

The same cycle also hardened the semantics of operator validity.  A local
response is now classified as \texttt{InvalidOperator} if any tangent column is
missing entirely.  One-sided fallback with a complete operator remains
\texttt{Degraded}; only a fully central, fully valid, non-regularised operator
is reported as \texttt{Ok}.  This change is small in code volume, but
scientifically important because it prevents an incomplete section operator
from being treated as merely ``low quality'' and then entering the FE\textsuperscript{2}
loop as if it were usable.

\section{Phase 2: Protocol Hardening for Macro--Micro Iteration}
\label{sec:ms-stab-phase2}

\subsection{Initial hypothesis}

The checkpointed FE\textsuperscript{2} loop added in the redesign was close to
publication-ready, but the failure semantics were still too optimistic and too
opaque for a release gate.

\subsection{Errors confirmed during audit}

Two issues were confirmed in the code:
\begin{enumerate}
  \item on a macro trial-step failure inside
    \texttt{IteratedTwoWayFE2}, the macro checkpoint was not restored
    explicitly before returning;
  \item \texttt{CouplingIterationReport} did not expose the termination reason,
    rollback execution, regularisation presence, or failed coupling sites.
\end{enumerate}

\subsection{Implemented change}

The orchestrator now records:
\begin{itemize}
  \item \texttt{CouplingTerminationReason};
  \item whether rollback was performed;
  \item whether relaxation changed the current iterate;
  \item whether regularisation was detected among accepted local responses;
  \item the set of hard-failure sites;
  \item a count of regularised submodels.
\end{itemize}

Most importantly, the macro checkpoint is now restored explicitly on macro-step
failure before the FE\textsuperscript{2} step returns.  This closes a real
semantic gap that could otherwise leave a dirty trial state in memory after a
reported failure.

\section{Phase 3: Physical Regression Matrix}
\label{sec:ms-stab-phase3}

\subsection{Added regression coverage}

The stabilization cycle introduced a small deterministic matrix of physical
reference cases using the real \texttt{NonlinearSubModelEvolver} path:
\begin{itemize}
  \item pure axial loading;
  \item pure bending about \(y\);
  \item pure bending about \(z\);
  \item pure shear along \(y\);
  \item pure shear along \(z\);
  \item pure torsion.
\end{itemize}

The tests check:
\begin{itemize}
  \item convergence of the local sub-model;
  \item availability of the homogenised response;
  \item positivity of the expected tangent diagonal;
  \item completion of the tangent diagnostics;
  \item use of \texttt{LinearizedCondensation} and the
    \texttt{forces\_consistent\_with\_tangent} flag in the elastic cases;
  \item dominance or mode-consistent prominence of the expected resultant;
  \item linearised consistency in the elastic axial case;
  \item boundary-reaction versus volume-average agreement in the elastic axial
    case;
  \item axial macro--micro power consistency.
\end{itemize}

\subsection{Observed numerical results}

The focused axial consistency test produced the following values on the local
execution used for this chapter:
\begin{align}
  \text{linearization\_error} &= 7.85096\times 10^{-5}, \\
  \text{operator\_gap} &= 2.84408\times 10^{-5}, \\
  \text{power\_gap} &= 2.84408\times 10^{-9}.
\end{align}
These are sufficiently small for the intended stabilization gate and, more
importantly, provide a reproducible numerical baseline for future refactors.

\subsection{Executed regression targets}

The following focused executables were rebuilt and executed locally after the
stabilization changes:
\begin{itemize}
  \item \texttt{fall\_n\_multiscale\_api\_test.exe}: 39 passed, 0 failed;
  \item \texttt{fall\_n\_evolver\_advanced\_test.exe}: 80 passed, 0 failed;
  \item \texttt{fall\_n\_heterogeneity\_test.exe}: 26 passed, 0 failed.
\end{itemize}

\section{Phase 4: Release Hardening and CI Surface}
\label{sec:ms-stab-phase4}

\subsection{Local reproducible gate}

The repository now contains
\texttt{scripts/ci\_multiscale\_stabilization.ps1}, which configures the build
directory if needed, builds the multiscale release gate, executes the key
regression tests, and compiles the PDF.  The script covers:
\begin{itemize}
  \item \texttt{fall\_n\_multiscale\_api\_test};
  \item \texttt{fall\_n\_micro\_solve\_executor\_test};
  \item \texttt{fall\_n\_evolver\_advanced\_test};
  \item \texttt{fall\_n\_beam\_test};
  \item \texttt{fall\_n\_steppable\_solver\_test};
  \item \texttt{fall\_n\_steppable\_dynamic\_test};
  \item \texttt{fall\_n\_lshaped\_multiscale};
  \item \texttt{fall\_n\_lshaped\_multiscale\_16};
  \item \texttt{fall\_n\_table\_multiscale};
  \item \texttt{fall\_n\_table\_cyclic\_validation};
  \item \texttt{latexmk -pdf}.
\end{itemize}

During this execution cycle, the script itself exposed two practical issues:
\begin{enumerate}
  \item the initial \texttt{param()} placement was invalid for PowerShell and
    had to be corrected;
  \item the document step required a defensive fallback to
    \texttt{pdflatex} on Windows when \texttt{latexmk} propagated a stale
    toolchain/runtime failure.
\end{enumerate}
After those corrections, the equivalent gate commands were executed explicitly
in the audit environment to verify the release surface.

For this follow-up execution cycle, the focused regression executables and the
\LaTeX{} document were revalidated again after the condensed tangent
implementation.  A fresh rebuild of the heavy example targets was also
attempted, but the harness hit its wall-clock timeout before completion; this
is therefore recorded as an incomplete integration check rather than silently
reported as green.

\subsection{Repository workflow}

The repository now also contains the workflow
\texttt{.github/workflows/multiscale-stability.yml}.  Its purpose is not to
claim that hosted CI is already production-proven, but to freeze the intended
release gate in repository form:
\begin{enumerate}
  \item configure an MSYS2 UCRT64 environment on Windows;
  \item install the numerical dependencies used by the project;
  \item configure CMake with Ninja;
  \item build the multiscale stability target set;
  \item run the focused regression executables;
  \item compile the \LaTeX{} document.
\end{enumerate}

\section{Phase 5: Residual Debt Ledger}
\label{sec:ms-stab-phase5}

The stabilization cycle deliberately left some debt explicit rather than
masking it under a ``green'' release summary.  Table
\ref{tab:ms-stab-residual-debt} records the items that remain open after this
execution.

\begin{table}[htb]
\centering
\small
\begin{tabular}{p{0.28\textwidth}p{0.14\textwidth}p{0.40\textwidth}}
\toprule
Residual item & Severity & Current disposition \\
\midrule
Dense condensed tangent still depends on explicit element-level local
linearisation support and keeps a finite-difference fallback &
high &
Documented explicitly; the primary path is now the condensed operator, but the
fallback remains necessary for unsupported local elements or singular free
blocks. \\
\texttt{NonlinearSubModelEvolver} still owns several responsibilities &
medium &
Improved in this cycle by extracting boundary-condition, state, and
homogenisation helpers; solver-core, diagnostics, and output concerns still
need separation. \\
MPI-distributed micro execution &
medium &
Communicator ownership is explicit, but the distributed execution engine is not
yet implemented. \\
Type-erased local-model path &
medium &
Left available only as experimental code until erased checkpoint/restore is
completed. \\
Legacy documentation warnings unrelated to this stabilization &
low/medium &
Inventoried below; not silently fixed in this cycle because they pre-existed
and are broader than the multiscale module. \\
\bottomrule
\end{tabular}
\caption{Residual debt after the stabilization execution.}
\label{tab:ms-stab-residual-debt}
\end{table}

\section{Documentation Warning Ledger}
\label{sec:ms-stab-doc-ledger}

The PDF compilation succeeded during this cycle, but the log still reports
pre-existing warnings outside the immediate stabilization work:
\begin{itemize}
  \item undefined references;
  \item multiply-defined labels;
  \item several hyperref string warnings associated with older chapters.
\end{itemize}

These warnings were intentionally not hidden or silently normalised.  The
stabilization decision was to keep the document compiling, add the current
chapter as an execution record, and leave the broader document cleanup as a
separate documentation-quality task.

\section{Closure of the Stabilization Cycle}
\label{sec:ms-stab-close}

The multiscale stabilization cycle v1 closed the most important gap between
``public API redesign'' and ``release-worthy behaviour'':
\begin{enumerate}
  \item the local tangent path now exposes its own uncertainty instead of
    silently zeroing columns;
  \item incomplete local section operators are now classified explicitly as
    invalid rather than being treated as weak-but-usable responses;
  \item the FE\textsuperscript{2} protocol reports why it terminated and
    whether it rolled back;
  \item the physical path now has deterministic axial, flexural, shear, and
    torsional regression coverage;
  \item the release gate exists both as a local script and as a repository
    workflow;
  \item the debt that remains is explicit, bounded, and technically
    understandable.
\end{enumerate}
